\subsection{Störeinflüsse bei Wechselstrommessbrücken}\label{sec:3-1-stoereinfluesse-wechselstrombruecken}

\subsubsection{Aufgabenstellung}

Es sind die Fehlerquellen der Wechselstrommessbrücken zu diskutieren. Und Unterdrückungen dieser anzugeben.

\subsubsection{Störeinflüsse}
Generell sind Wechselstrombrücken empfindlich gegenüber externen Feldern sowie der
Qualität der Speisespannung. Im Folgenden werden die wichtigsten Einflüsse samt
konkreter Maßnahmen kurz zusammengefasst.

\paragraph{Magnetische Felder}
Zeitlich veränderliche magnetische Felder induzieren Schleifenspannungen \cite[S.~49, Gl.~2.24]{EMTPraktikum2025}
\[
	u_{\text{ind}} \propto A\,\frac{\mathrm{d} B}{\mathrm{d}t}
\]
und verschieben den Brückenabgleich. \textit{Maßnahmen:} kleine Schleifenflächen,
verdrillte/bifilare Leitungen zwischen Brückeneckpunkten, kurze Abstände der Knoten,
magnetisch ruhige Aufstellung, ggf.\ Ferritkerne für kleinere Streufelder.

\paragraph{Elektrische Felder \& parasitäre Kapazitäten}
Potenzialunterschiede zwischen Brücken-knoten und Erde bilden parasitäre Kapazitäten.
Ströme über diese erzeugen zusätzliche Spannungsabfälle an den
Brückenimpedanzen.

\textit{Maßnahmen:} Eine Schirmung der Leitungen gegen Erde reduziert den
Einfluss externer elektrischer Felder. Die Positionierung des Erdknotens ist
relevant, oft wird ein Kontakt des Nullinstruments geerdet. \cite[S.50]{EMTPraktikum2025}

\paragraph{Versorgung}
Oberwellen und Frequenzdrift der Quelle sowie Sättigungseffekte des
Übertragers (z.B. Übersetzung $4.7:1$) erzeugen hochfrequente Anteile und
Phasenfehler. Dieser Einfluss erschwert den Abgleich, da kein reines Sinussignal mehr
vorliegt.

\textit{Maßnahmen:} Zunächst ist die Qualität der Quelle zu beachten. Die
Quelle soll die Eigenschaften: wenig Oberwellen, frequenzstabile Quelle und
ausreichend große Sättigungsreserve erfüllen.
Die Quelle soll so eingestellt werden, dass weder niedrige Frequenzen (z.B. Netzfrequenz)
noch hohe Frequenzen wo die Kapazitäten stärkeren einfluss haben, einen großen einfluss auf das
Brückensignal haben. \cite[S.50]{EMTPraktikum2025}
